% !TeX root = main_long.tex
\section{Results}
\label{sec:coverage}

\subsection{RQ1: coverage of VCF constructs and meaning}
\label{sec:evidence}

\paragraph{Declarations and assessment corpus.}
All \textbf{333} explicit reserved Number/Type declarations extracted from VCF 4.1--4.5 match the version registries, with no missing or differing row.
VCF 4.5's registry contains 122 reserved definitions: 51 INFO and 71 FORMAT, including three parameterized base-modification families.
These are different populations: the comparison includes only explicit, comparable declarations, while older specifications also define reserved keys in prose without complete Number/Type information.
Agreement is therefore weighted towards VCF 4.3--4.5 and does not settle declaration--prose inconsistencies, such as CNL/CNP declaring \texttt{Number=G} while describing copy-number ordering.

The source audit yielded \textbf{94 requirements}, \textbf{38 fixtures}, \textbf{210 cases}, and \textbf{79 distinct queries}, linked to specification passages and expected answers.
Nine fixtures reproduce published specification examples exactly; the others are minimal files instantiating stated rules.
Availability of defensible examples influenced test selection, leaving the demonstrated subset uneven across the specification.
Interpretation statements also vary in specificity: 56 are requirement-specific, 29 generic, and nine explicitly unassessed.
During our review we identified inconsistent specification examples, including three issues (such as \href{https://github.com/samtools/hts-specs/issues/869}{hts-specs\#869}\footnote{\url{https://github.com/samtools/hts-specs/issues/869}}) reported to the maintainers and one previously reported issue, documented in the experimental repository.
Where an interpretation could be established, affected requirements were scored; unresolved cases remained unassessed and counted against coverage.

\paragraph{Executed coverage.}
The 210 cases exercise \textbf{53 of 94 requirements} in at least one applicable version; 41 have no registered test.
All 210 preservation tests passed in each profile, as did all 210 expanded-structure tests.
Condensed structure passed 165 and failed 45 tests, all concerning sample-level FORMAT access held in vectors rather than individual resources.
Every passing case rejected the empty-graph control and changed its answer under at least one referenced-predicate deletion.
Table~\ref{tab:requirement-results} reports requirements rather than cases, separately by applicable VCF version.

\begin{table}[htbp]
\centering
\caption{
  Demonstrated requirements by VCF version, sample profile, and retrieval axis. 
  Each Demonstrated column assesses the same $N$ applicable requirements: demonstrated $+$ ND $+$ U $=N$, with percentages calculated over $N$. 
  U denotes untested requirements; ND denotes tested failures and is nonzero only for condensed structure ($^\dagger$). 
  No requirement was partially demonstrated.
}
\label{tab:requirement-results}
\small
\begin{tabular*}{\linewidth}{@{\extracolsep{\fill}}lrrrrrrr@{}}
\toprule
 & & \multicolumn{4}{c}{Demonstrated} & & \\
\cmidrule(lr){3-6}
 & & \multicolumn{2}{c}{\emph{preservation}} & \multicolumn{2}{c}{\emph{structure}} & & \\
\cmidrule(lr){3-4} \cmidrule(lr){5-6}
VCF & $N$ & \textit{expanded} & \textit{condensed} & \textit{expanded} & \textit{condensed} & ND$^{\dagger}$ & U \\
\midrule
4.1 & 69 & 32 (46.4\%) & 32 (46.4\%) & 32 (46.4\%) & 26 (37.7\%) &  6 & 37 \\
4.2 & 72 & 34 (47.2\%) & 34 (47.2\%) & 34 (47.2\%) & 28 (38.9\%) &  6 & 38 \\
4.3 & 74 & 35 (47.3\%) & 35 (47.3\%) & 35 (47.3\%) & 29 (39.2\%) &  6 & 39 \\
4.4 & 83 & 38 (45.8\%) & 38 (45.8\%) & 38 (45.8\%) & 31 (37.3\%) &  7 & 45 \\
4.5 & 91 & 51 (56.0\%) & 51 (56.0\%) & 51 (56.0\%) & 37 (40.7\%) & 14 & 40 \\
\bottomrule
\end{tabular*}
\end{table}

For VCF 4.5, 51 of 91 requirements (56.0\%) are demonstrated in the expanded profile.
Of the 40 unassessed requirements, 37 await fixtures and queries, one requires further interpretation, and two require coverage of all reserved INFO and FORMAT keys.
Thus, most missing coverage reflects absent tests; the tested structural failures concern the condensed sample representation.
The corresponding expanded tests pass, demonstrating that those values can be expressed structurally within VCF Core.

\paragraph{Curated inventory.}
All \textbf{104} inventoried VCF 4.5 logical constructs were classified as preserved and structured; \textbf{87} additionally had a rule enforcing a stated requirement~\cite{vcfc210}.
Enforcement was less complete for headers, lexical rules, genotypes, and structural variation.
Evidence included queries, regression probes, shapes or decoded checks, and materialized fixtures.
These are curated representation judgments supported by different evidence types, rather than 104 independently executed retrieval tests.

\subsection{RQ2: behaviour across producers and profiles}
\label{sec:crossproducer}

Re-executing the 210 cases on two axes and two profiles gives \textbf{840 checks}: \textbf{693 pass for both producers}, \textbf{45 fail for both}, and \textbf{102 differ} across 15 requirements.
Every disagreement is a materializer pass and converter failure; none shows the inverse.
The 45 shared failures are condensed-structure queries requiring per-sample resources.
Of the 53 exercised requirements, 38 agree on every check and 32 pass on every check.
Inspection of each disagreement identifies the causes in Table~\ref{tab:divergence}.
No observed divergence requires a representational mechanism absent from the pinned vocabulary; every missing property is already declared.

\begin{table}[htbp]
\caption{Causes of disagreement across 15 requirements.}
\label{tab:divergence}
\centering\small
\begin{tabularx}{\linewidth}{@{}>{\raggedright\arraybackslash}p{0.29\linewidth}c>{\raggedright\arraybackslash}X@{}}
\toprule
Cause & Reqs. & Finding \\
\midrule
Unemitted properties & 12 & Twenty declared properties absent, including accessors for gVCF blocks, breakend mates, and allele indices. \\
Incomplete decomposition & 1 & Tandem-repeat summary emitted without per-allele components. \\
Permitted datatype difference & 1 & QUAL uses \term{vcfc:VCFFloat} or \term{xsd:decimal}; both satisfy the shape but differ under \term{DATATYPE()}. \\
Encoding convention & 1 & Different handling of an empty local-allele list (R27). \\
\bottomrule
\end{tabularx}
\end{table}

An earlier comparison also exposed a silent local-allele error missed by the converter's existing tests.
It associated \term{Number=LA} and \term{Number=LR} values with record alleles positionally instead of resolving the sample's local subset.
Under \texttt{LAA=2,4}, depths \texttt{20,30,10} belong to global alleles 0, 2, and 4, but were attached to 0, 1, and 2.
That defect and eight further divergences were repaired~\cite{vcfrConverter2026}.
In the evaluated release, all local-allele requirements pass for both producers on expanded-profile checks; the condensed-profile R27 divergence remains.

\paragraph{Semantic round-trip and profile comparison.}
Across 38 fixtures, reconstruction from structured properties recovers the seven fixed columns of all \textbf{109 records}, all \textbf{151 sample genotypes}, all \textbf{431 header lines}, and \textbf{400 of 401} INFO entries and non-GT FORMAT subfields.
The exception is an empty \texttt{LAA} list: the materializer writes an empty value, while the converter emits the field without a value triple.
This loses the distinction between a present, empty list, meaning reference-only local alleles, and a field without a value.
The round-trip therefore reaches R27 through a complementary check.
The associated \texttt{LAD} item retains its \term{forAllele} link, so the error does not remove that depth association.

\label{sec:profilesresult}
The sample profiles agree on \textbf{352 of 420} case-and-axis comparisons.
All 68 disagreements concern FORMAT values or genotypes within samples, spanning 15 requirements.
Questions about files, headers, records, and alleles agree throughout.
The profiles therefore agree above the sample level, while the condensed vectors require decoding that direct structural queries do not perform.

\subsection{RQ3: integration with retained source context}
\label{sec:integration}

The query returns exactly the two expected alterations: \texttt{chr1:1~G>C} at depth 30 and \texttt{chr1:2~A>T} at depth 25.
Each is traced to its file, record, sample, and the FORMAT declaration defining \texttt{LAD}'s cardinality, and every row identifies the file actually converted.
All three expected exclusions are confirmed individually: \texttt{chr1:1~G>A}, \texttt{chr1:2~A>C}, and \texttt{chr1:3~C>G} are outside the respective local-allele subsets; the last would also fail the depth threshold under the positional reading.

A positional interpretation of \texttt{Number=LR} also returns two rows with provenance, but associates the depths with \texttt{chr1:1~G>A} and \texttt{chr1:2~A>C} instead.
Thus, plausible values, the expected row count, and provenance can coexist with incorrect allele associations.
Two exclusions change only when the sample-specific \texttt{LAA} mapping is resolved.
This counterexample was defined from the fixture and specification rule before query execution; it illustrates the consequence of the rule, rather than measuring an error in a particular converter.
The result demonstrates the utility of local-allele-aware representation for this synthetic example, with no claim of clinical relevance or correctness across all VCF data.

\subsection{RQ4: relation to existing VCF semantic models}
\label{sec:models}

Table~\ref{tab:models} summarizes the inspected artifacts.
HERO-Genomics and GVO retain INFO as one literal through \term{hero:vcfInfo} and \term{gvo:info}, requiring consumers to parse the text and supply interpretation rules~\cite{cazzaro2025hero,kawashima2023gvo}.
GFVO directly models field meanings such as \term{AlleleCount} and \term{MappingQuality}, but leaves the declaration mechanism outside its scope~\cite{baran2015gfvo}.
VCF2RDF publishes declaration terms, including \term{Number}, \term{Type}, and \term{Description}, but no allele-index term, so allele association requires additional conventions~\cite{penha2017vcf2rdf}.
GVO's predominantly S verdicts reflect its focus on classifying biological variation rather than representing source-file relationships.

\begin{table}[htbp]
\caption{Support in inspected artifacts: E, explicitly represented; C, requiring additional conventions; S, outside stated scope. These are inspection judgments about VCF interpretation context.}
\label{tab:models}
\centering\small
\begin{tabularx}{\linewidth}{@{}Xccccc@{}}
\toprule
Can a consumer\ldots & VCF Core & GFVO & HERO & GVO & VCF2RDF \\
\midrule
F1\quad trace a record to its file & E & E & E & S & C \\
F2\quad read declaration ID, Number, and Type & E & S & S & S & E \\
F3\quad read the declared VCF version & E & S & S & S & C \\
F4\quad address ALT allele $n$ by index & E & C & C & S & C \\
F5\quad separate ploidy, allele calls, and phasing & E & C & C & S & C \\
F6\quad get a (sample, field) value unsplit & E & C & C & S & C \\
F7\quad bind a \texttt{Number=A/R/G} item to its allele & E & C & C & S & C \\
F8\quad distinguish explicit \texttt{.} from unstated & E & S & C & C & C \\
F9\quad identify the sample encoding & E & S & S & S & S \\
\bottomrule
\end{tabularx}
\end{table}

VCF Core's distinguishing role is the explicit representation of source interpretation context: declarations, version rules, allele indices, and sample scope.
These findings concern the published artifacts and do not show that other models cannot support the same information through extensions.
Separately, VCF Core provides \textbf{123 curated alignments} \footnote{\url{https://github.com/ecrum19/vcf-core-vocabulary/tree/main/mappings}}, recorded in SSSOM with SKOS predicates: 57 exact, 40 close, 16 related, and ten narrow matches~\cite{matentzoglu2022sssom,skos2009}.
These graded term mappings support connections to other models without establishing identity between instances.
